\documentclass[11pt]{article}
\usepackage[T1]{fontenc}
\usepackage{lmodern}
\usepackage{amsmath,amssymb,amsthm,mathtools}
\usepackage{microtype}
\usepackage[margin=1in]{geometry}
\usepackage[hidelinks]{hyperref}

\newtheorem{theorem}{Theorem}[section]
\newtheorem{lemma}[theorem]{Lemma}
\newtheorem{proposition}[theorem]{Proposition}
\newtheorem{corollary}[theorem]{Corollary}
\theoremstyle{definition}
\newtheorem{remark}[theorem]{Remark}

\newcommand{\N}{\mathbb N}
\newcommand{\R}{\mathbb R}
\newcommand{\1}{\mathbf 1}
\newcommand{\Om}{\Omega}
\newcommand{\rhoL}{\rho_L}
\newcommand{\Pmin}{P^-}
\newcommand{\floor}[1]{\left\lfloor #1\right\rfloor}
\newcommand{\abs}[1]{\left|#1\right|}
\newcommand{\set}[1]{\left\{#1\right\}}
\newcommand{\cref}[1]{\ref{#1}}
\newcommand{\llp}[1]{\mathrel{\ll_{#1}}}
\DeclareMathOperator{\rad}{rad}
\newcommand{\Pplus}{P^+}
\newcommand{\adm}{\mathrm{adm}}
\newcommand{\Z}{\mathbb Z}


\title{Admissibility variants of Erd\H{o}s Problem 327:\\
the multipliers \(k\geq2\)}
\author{Donald Della Pietra}
\date{}
\hypersetup{
  pdftitle={Admissibility variants of Erdos Problem 327: the multipliers k >= 2},
  pdfauthor={Donald Della Pietra},
  pdfsubject={Unit fractions and Erdos Problem 327},
  pdfkeywords={unit fractions, multiplicative functions, anatomy of integers, Erdos Problem 327}
}

\begin{document}
\maketitle

\begin{abstract}
For \(k\geq1\) call a set \(A\) \(k\)-admissible if \(a+b\nmid kab\) for all
distinct \(a,b\in A\), and let \(f_k(N)\) be the largest size of a
\(k\)-admissible subset of \(\{1,\ldots,N\}\). The case \(k=1\) is the first
question of Erd\H{o}s Problem 327 and \(k=2\) is the second; Sawin proved
\(f_2(N)\gg N\), and the companion manuscript \cite{DellaPietra2026} proved
\(f_1(N)\geq(1/2+\varepsilon)N\).

We first settle what the correct benchmark is for each \(k\). A periodic set
is \(k\)-admissible if and only if every one of its residue classes consists
of odd numbers and \(k\) is odd. Hence for odd \(k\) the odd numbers are
\(k\)-admissible and are the densest periodic \(k\)-admissible set, while for
even \(k\) \emph{no} nonempty periodic set is \(k\)-admissible. The benchmark
is therefore \(1/2\) for odd \(k\) and \(0\) for even \(k\); this is the
analogue, for the present problem, of a parity trap, and it changes which
target is meaningful.

We then show that the machinery of \cite{DellaPietra2026} transfers with
\(k\)-dependent constants and an unchanged parameter window, provided its
dilation \(2\) is replaced by \(\lambda(k)\), the least power of two that is at
least \(k\). That replacement is necessary, not cosmetic: with the dilation
\(2\) the construction fails outright whenever \(3\mid k\) or \(5\mid k\),
because \(\{n,2n\}\) and \(\{2n,3n\}\) are then never \(k\)-admissible. We
obtain:
\begin{itemize}
\item for every odd \(k\) an absolute \(\varepsilon_k>0\) with
      \(f_k(N)\geq(\tfrac12+\varepsilon_k)N\) for large \(N\)
      (Theorem~\ref{thm:odd-main});
\item for every \(k\geq1\) an absolute \(c_k>0\) with \(f_k(N)\geq c_kN\)
      (Theorem~\ref{thm:all-k}). This is new for even \(k\geq4\); at \(k=2\)
      it recovers Sawin's theorem, and for odd \(k\) it is superseded by the
      previous item and, indeed, already by the odd numbers.
\end{itemize}
We do \emph{not} obtain \(f_2(N)\geq(1/2+\varepsilon)N\), the question raised
in the Problem 327 discussion thread, and Section~\ref{sec:even} isolates the exact
obstruction: for even \(k\) the method has no dense automatically admissible
layer to start from, and the dilation identity shows that no rescaling
produces one. All numerical claims are certified by exact rational arithmetic
in an accompanying standalone verifier. A Lean formalization is planned and
is not part of this manuscript.
\end{abstract}

\section{Introduction}\label{sec:intro}

For \(k\geq1\), call a set \(A\) of positive integers \(k\)-admissible if
\begin{equation}\label{eq:kadm}
 a+b\nmid kab\qquad(a,b\in A,\ a\ne b),
\end{equation}
and put
\[
 f_k(N)=\max\{\abs{A}:A\subseteq\{1,\ldots,N\},\ A\text{ is }k\text{-admissible}\}.
\]
Erd\H{o}s and Graham \cite[p.~37]{ErdosGraham1980} asked whether a
\(1\)-admissible set can be substantially more than the odd numbers, and
whether every \(2\)-admissible set has size \(o(N)\); these are the two
questions of Problem 327 \cite{ErdosProblem327}. In unit-fraction language
\(a+b\mid ab\) says that \(1/a+1/b\) is a unit fraction, and \(a+b\mid2ab\)
says that the average of \(1/a\) and \(1/b\) is a unit fraction.

The state of the art we start from is the following.
\begin{itemize}
\item Sawin \cite{Sawin2026} proved \(f_2(N)\gg N\), answering the second
      question negatively. The constant is not made explicit.
\item The companion manuscript \cite{DellaPietra2026} proved
      \(f_1(N)\geq(1/2+\varepsilon)N\) for an absolute \(\varepsilon>0\).
\item In the discussion thread \cite{Erdos327Thread}, leon2k2k2k obtained the
      upper bounds \(f_1(N)\leq(0.7769+o(1))N\),
      \(f_2(N)\leq(0.7630+o(1))N\) and \(f_3(N)\leq(0.6089+o(1))N\)
      \cite{Leon2026}, and asked whether
      \(f_2(N)\geq(1/2+o(1))N\).
\end{itemize}
For odd \(k\) the odd numbers give the trivial bound
\(f_k(N)\geq N/2+O(1)\), as we recall in Section~\ref{sec:benchmark}. Beyond
that, no lower bound for \(f_k\) with \(k\geq3\) --- in particular none at all
for even \(k\geq4\), where the trivial bound is unavailable --- appears in the
literature or in the thread; a dated search is reported in
Section~\ref{sec:repro}.

\subsection{The correct benchmark}

Before asking whether one can beat density \(1/2\) it is necessary to know
whether \(1/2\) is available at all. It is not, for even \(k\). The following
is proved in Section~\ref{sec:benchmark}; it is sharp, and it is what makes the
odd and even cases genuinely different problems.

\begin{theorem}[periodic benchmark]\label{thm:benchmark}
Let \(k,M\geq1\) and \(0\leq r<M\), and let
\(C_{r,M}=\{n\geq1:n\equiv r\bmod M\}\).
\begin{enumerate}
\item If \(k\) is odd, \(M\) is even and \(r\) is odd, then \(C_{r,M}\) is
      \(k\)-admissible.
\item In every other case \(C_{r,M}\) contains distinct \(a,b\) with
      \(a+b\mid kab\); an explicit such pair is displayed in
      \eqref{eq:witness}.
\end{enumerate}
Consequently, a union of residue classes to a common modulus is
\(k\)-admissible if and only if each of its classes is of type (1). For odd
\(k\) every periodic \(k\)-admissible set consists of odd numbers and so has
density at most \(1/2\), with equality for the odd numbers; for even \(k\)
the only periodic \(k\)-admissible set is the empty set.
\end{theorem}

So the benchmark density is \(1/2\) for odd \(k\) and \(0\) for even \(k\),
and ``beating the odd numbers'' is a meaningful goal exactly when \(k\) is
odd. For even \(k\) even a single residue class must be pruned, and every
construction has to be aperiodic from the outset.

\subsection{Results}

\begin{theorem}[odd multipliers]\label{thm:odd-main}
For every odd \(k\geq1\) there are constants \(\varepsilon_k>0\) and \(N_k\)
such that
\[
 f_k(N)\geq\left(\frac12+\varepsilon_k\right)N
 \qquad(N\geq N_k).
\]
\end{theorem}

For \(k=1\) this is \cite[Theorem 1.1]{DellaPietra2026} and contains nothing
new; the content here is that the argument transfers to every odd \(k\) with
the \emph{same} certified parameter window and only \(k\)-dependent
multiplicative constants. Combined with the upper bound of
\cite{Leon2026} --- itself a thread posting, independently checked but not
refereed --- this would give, for example,
\[
 \tfrac12<\liminf_{N\to\infty}\frac{f_3(N)}{N}
 \leq\limsup_{N\to\infty}\frac{f_3(N)}{N}\leq0.6089 .
\]

\begin{theorem}[all multipliers]\label{thm:all-k}
For every \(k\geq1\) there is a constant \(c_k>0\) such that
\(f_k(N)\geq c_kN\) for all sufficiently large \(N\).
\end{theorem}

For odd \(k\), Theorem~\ref{thm:all-k} is trivial: the odd numbers already
have density \(1/2\). Its content is for even \(k\), where by
Theorem~\ref{thm:benchmark} no such trivial set exists. At \(k=2\) it reproves
Sawin's theorem by the route of \cite[\S4]{DellaPietra2026}; for even
\(k\geq4\) it is, to our knowledge, the first lower bound of any kind. It uses
only the source half of the machinery, not the two-endpoint estimate, and is
therefore the more robust of the two results. Because both
analytic inputs carry unspecified absolute constants, \(c_k\) is not
effective as stated; see Section~\ref{sec:gaps}.

\subsection{What does not close}

We do not prove \(f_2(N)\geq(1/2+\varepsilon)N\), nor any bound for even \(k\)
beyond positive density. Section~\ref{sec:even} records three separate reasons,
each of which is proved rather than asserted: the absence of a periodic
layer (Theorem~\ref{thm:benchmark}), the dilation identity
(Proposition~\ref{prop:dilation}), which shows that every rescaled layer of an even-\(k\)
construction is again an even-multiplier problem, and the fact that the
roughness cutoff, which is what buys the saving in the source estimate, caps
the density of the source at \(O(\rhoL)\). We also record, in
Proposition~\ref{prop:fibration}, a corrected form of the largest-prime-factor
observation from the thread, and show that its contraction constant
\(\log2<1\) prevents it from bootstrapping any density.

\subsection{Method and organisation}

Section~\ref{sec:benchmark} is self-contained and elementary. Section~\ref{sec:inputs}
recalls the two analytic inputs and the centered anatomy of
\cite{DellaPietra2026} verbatim; nothing there is new.
Section~\ref{sec:source} generalises the source estimate from the modulus \(2\) to an
arbitrary fixed modulus \(m\) and proves Theorem~\ref{thm:all-k}.
Section~\ref{sec:mixed} generalises the two-endpoint mixed estimate to odd
\(k\), and this is where the dilation \(\lambda(k)\) of \eqref{eq:lambda} is
introduced and where Remark~\ref{rem:lambda-two} explains why the naive
dilation \(2\) is fatal. Section~\ref{sec:assembly} assembles
Theorem~\ref{thm:odd-main}. Section~\ref{sec:even} treats the
even case and the obstruction. Section~\ref{sec:gaps} lists every step that is
imported, flagged, or unverified, and Section~\ref{sec:repro} describes the
verifier, the prior-work search and the AI-use disclosure.

\section{The correct benchmark}\label{sec:benchmark}

Everything in this section is elementary and is confirmed exhaustively by the
verifier.

\begin{lemma}[conflict criterion]\label{lem:criterion}
Let \(a\ne b\) be positive integers, \(g=(a,b)\), and \(x=(a+b)/g\). Then for
every \(k\geq1\),
\[
 a+b\mid kab
 \qquad\Longleftrightarrow\qquad
 \frac{x}{(x,k)}\ \Big|\ g .
\]
\end{lemma}

\begin{proof}
Write \(a=gu\), \(b=gv\) with \((u,v)=1\), so \(x=u+v\) and \(a+b=gx\),
\(kab=kg^2uv\). Thus \(a+b\mid kab\) if and only if \(x\mid kguv\). Now
\((x,u)=(u+v,u)=(v,u)=1\) and likewise \((x,v)=1\), so \((x,uv)=1\) and the
condition is \(x\mid kg\). Finally, for every prime \(p\),
\[
 v_p(x)\leq v_p(k)+v_p(g)
 \iff
 v_p(g)\geq v_p(x)-\min\bigl(v_p(x),v_p(k)\bigr)
 =v_p\!\left(\frac{x}{(x,k)}\right),
\]
which is the stated divisibility.
\end{proof}

\begin{corollary}[two-adic restriction]\label{cor:twoadic}
Let \(a,b\) be distinct odd integers with \(a+b\mid kab\). Then
\(v_2(a+b)\leq v_2(k)\). In particular, for \(k=2\) one has
\(a\equiv b\pmod4\): the conflict graph of the multiplier \(2\) on the odd
numbers has no edge between the classes \(1\) and \(3\) modulo \(4\).
\end{corollary}

\begin{proof}
Here \(g\mid a\) is odd, so \(v_2(x)=v_2(gx)=v_2(a+b)\). By
Lemma~\ref{lem:criterion} we have \(x\mid kg\), whence
\(v_2(x)\leq v_2(k)\). For \(k=2\) this gives \(v_2(a+b)\leq1\), and since
\(a,b\) are odd, \(a+b\equiv2\pmod4\), i.e.\ \(a\equiv b\pmod4\).
\end{proof}

\begin{proof}[Proof of Theorem~\ref{thm:benchmark}]
(1) If \(a,b\) are odd then \(a+b\) is even while \(kab\) is odd, so
\eqref{eq:kadm} holds. A class \(C_{r,M}\) with \(M\) even and \(r\) odd
consists of odd numbers.

(2) Suppose we are not in case (1), i.e.\ \(k\) is even, or \(M\) is odd, or
\(r\) is even. Put
\[
 u=1,\qquad v=1+2M,\qquad x=u+v=2(1+M),\qquad y=\frac{x}{(x,k)} .
\]
We claim \((y,M)\mid r\). Let \(p\) be a prime dividing \(M\). If \(p\) is
odd then \(x=2+2M\equiv2\pmod p\), so \(p\nmid x\) and hence \(p\nmid y\).
If \(p=2\), then \(M\) is even, so \(1+M\) is odd and \(v_2(x)=1\); therefore
\(v_2(y)=1-\min(1,v_2(k))\), which is \(0\) when \(k\) is even and \(1\) when
\(k\) is odd. Consequently
\[
 (y,M)=
 \begin{cases}
 2,&\text{\(M\) even and \(k\) odd},\\
 1,&\text{otherwise},
 \end{cases}
\]
and in each of the three cases excluded from (1) we get \((y,M)\mid r\):
if \(k\) is even or \(M\) is odd then \((y,M)=1\), and if \(M\) is even,
\(k\) is odd and \(r\) is even then \((y,M)=2\mid r\).

Since \((y,M)\mid r\), the simultaneous congruences \(g\equiv0\pmod y\) and
\(g\equiv r\pmod M\) have a positive solution \(g\). Put
\begin{equation}\label{eq:witness}
 a=g,\qquad b=g(1+2M).
\end{equation}
Then \(a\equiv r\) and \(b\equiv g\equiv r\pmod M\), and \(a\ne b\) because
\(1+2M>1\). Moreover \((a,b)=g\) and \((a+b)/g=1+v=x\), so
Lemma~\ref{lem:criterion} applies with this \(g\) and this \(x\): since \(y\mid g\)
we conclude \(a+b\mid kab\).

For the final statement, a union \(A\) of classes modulo \(M\) is
\(k\)-admissible only if each individual class is, since a conflict inside one
class is a conflict inside \(A\); and conversely, if every class of \(A\) is
of type (1) then \(A\) consists of odd numbers and is \(k\)-admissible by (1).
For odd \(k\) the admissible classes are exactly the odd classes, whose union
has density at most \(1/2\); for even \(k\) case (1) is vacuous.
\end{proof}

\begin{remark}
The threshold in Theorem~\ref{thm:benchmark} is delicate. The class
\(7\bmod10\) contains no \(2\)-conflict among its elements below \(897\); its
first three conflicts are \((117,897)\), \((357,1377)\) and \((77,1617)\). A
purely computational search in a box of size a few hundred would have reported
the wrong benchmark. Note also that the witness \eqref{eq:witness} is explicit
but not minimal: at \(k=2\), \(M=10\), \(r=7\) it has \(y=11\) and \(g=77\),
returning the third of those pairs and not the first. The construction is
designed to be uniform in \((k,M,r)\), not sharp.
\end{remark}

\begin{proposition}[dilation identity]\label{prop:dilation}
Let \(\lambda\) be a positive integer. A set \(A\) satisfies
\(\lambda a+\lambda b\nmid k(\lambda a)(\lambda b)\) for all distinct
\(a,b\in A\) if and only if \(A\) is \(k\lambda\)-admissible. Consequently
\[
 f_k(N)\geq f_{k\lambda}\!\left(\floor{N/\lambda}\right).
\]
\end{proposition}

\begin{proof}
\(\lambda(a+b)\mid k\lambda^2ab\) if and only if \(a+b\mid k\lambda ab\).
\end{proof}

This is the identity behind the two-layer construction: the even part of a
\(1\)-admissible set is \(2\) times a \(2\)-admissible set, an equivalence
already recorded in the thread. It also shows that for even \(k\) every
dilated layer carries an even multiplier, so Theorem~\ref{thm:benchmark} denies a
periodic layer at every scale.

\section{Notation and imported analytic inputs}\label{sec:inputs}

All logarithms are natural, \(\Om(n)\) is the number of prime factors of
\(n\) with multiplicity, \(\Pmin(n)\) is the least prime factor of \(n\) with
\(\Pmin(1)=+\infty\), and for real \(L>2\)
\[
 \Om_L(n,x)=\sum_{\substack{p^\nu\parallel n\\L\leq p\leq x}}\nu,
 \qquad
 \rhoL=\prod_{p<L}\left(1-\frac1p\right).
\]
An integer is \(L\)-rough if it has no prime factor below \(L\), and
\((L,A,K)\)-regular if
\begin{equation}\label{eq:regular}
 \Om_L(n,x)\leq A\log\frac{\log x}{\log L}+K\qquad(x\geq L).
\end{equation}
For an integer \(n\) and a cutoff \(L\) we write \(n=n_{<L}n_{\geq L}\) for
the factorisation into the \(L\)-smooth and \(L\)-rough parts.

We import three statements from \cite{DellaPietra2026} without change. The
first two are the analytic inputs; the third is the centered tail estimate,
whose proof involves no multiplier and is reproduced there in full.

\begin{proposition}[one-variable mean value; {\cite[Prop.~2.1]{DellaPietra2026}},
after {\cite[Thm.~III.3.5]{Tenenbaum2015}}]\label{prop:one-variable}
Let \(g\) be a nonnegative multiplicative function with
\(\sum_{p\leq y}g(p)\log p\leq C_1y\) for all \(y\geq1\) and
\(\sum_p\sum_{\nu\geq2}g(p^\nu)\log(p^\nu)p^{-\nu}\leq C_2\). Then
\[
 \sum_{n\leq Y}g(n)\ll(C_1+C_2+1)Y
 \prod_{p\leq Y}\left(1-\frac1p+\sum_{\nu\geq1}\frac{g(p^\nu)}{p^\nu}\right)
\]
with an absolute implied constant.
\end{proposition}

\begin{proposition}[three linear forms; {\cite[\S2]{DellaPietra2026}},
after {\cite[Thm.~3.1]{BretecheTenenbaum2026}}]\label{prop:dBT}
Let \(Q_1,Q_2,Q_3\in\Z[U,V]\) be fixed, primitive, pairwise nonproportional
linear forms, and let \(F:\N^3\to\R_{\geq0}\) be multiplicative in the sense
that \(F(a_1b_1,a_2b_2,a_3b_3)\leq F(b_1,b_2,b_3)\) whenever
\((a_1a_2a_3,b_1b_2b_3)=1\). Then, on every box of side lengths comparable to
\(X\),
\begin{equation}\label{eq:dBT-shape}
 \sum_{(u,v)}F\bigl(Q_1(u,v),Q_2(u,v),Q_3(u,v)\bigr)
 \ll X^2E_{F,Q}(CX)\prod_{3<p\leq CX}\left(1-\frac{\rho_Q^+(p)}{p^2}\right),
\end{equation}
where \(\rho_Q^+(p)=\#\{(u,v)\bmod p:p\mid Q_1Q_2Q_3(u,v)\}\), \(C\) depends
only on the forms and the box aspect ratio, and \(E_{F,Q}\) is bounded by the
product of its local prime-power factors. For the triples used below a local
pattern with exactly one positive valuation \(e\) has density
\begin{equation}\label{eq:axis-density}
 \frac{(1-p^{-1})^2}{p^e}
\end{equation}
at every odd prime, and patterns with two or three positive valuations
vanish. The implied constant is uniform over the class once the forms and
the box aspect ratio are fixed; this is what permits \(F\) to depend on the
cutoff \(L\).
\end{proposition}

\begin{lemma}[centered maximal tail;
{\cite[Lem.~3.1]{DellaPietra2026}}]\label{lem:centered-tail}
Fix \(1<z<2\) and \(A>0\) with \(A\log z>z-1\). Uniformly for real \(L>2\),
\(Y\geq1\) and \(K\geq0\),
\[
 \#\{n\leq Y:n\text{ is not }(L,A,K)\text{-regular}\}\ll_{A,z}Yz^{-K},
\]
and among \(L\)-rough integers the same count is
\(\ll_{A,z}Y\rhoL z^{-K}\).
\end{lemma}

\section{The source estimate for an arbitrary modulus}\label{sec:source}

Fix an integer \(m\geq1\). This section proves the existence of an
\(m\)-admissible set of positive density. The argument is
\cite[\S4]{DellaPietra2026} with \(2\) replaced by \(m\); we display the four
places where the multiplier intervenes, each of which costs only a constant.

Throughout, \(\delta\in(0,1)\), \(1<A<3/(2\log4)\), \(K\geq0\) and
\(L>m\) are fixed, and
\[
 S=S(N;m,L,A,K)=\{b\in[\delta N,N]:\Pmin(b)\geq L,\
 b\text{ is }(L,A,K)\text{-regular}\}.
\]

\begin{lemma}[source density;
{\cite[Lem.~4.1]{DellaPietra2026}}]\label{lem:source-density}
There is \(K_0=K_0(A,\delta)\), independent of \(L\) and of \(m\), such that
for \(K\geq K_0\), every fixed \(L>2\) and all large \(N\),
\(\abs{S}\geq\frac{1-\delta}{2}N\rhoL\).
\end{lemma}

Call \(b\in S\) \emph{bad} if there is \(c\leq N\), \(c\ne b\), with
\begin{equation}\label{eq:source-conflict}
 b+c\mid mbc,\qquad \Om(c)\leq\Om(b),
\end{equation}
and let \(T=T(N;m,L,A,K)\) be the number of bad \(b\).

\begin{lemma}[generalised Sawin coordinates]\label{lem:source-coords}
Let \(b\in S\) and \(c\) satisfy \eqref{eq:source-conflict}. Put
\(g=(b,c)\), \(u=b/g\), \(v=c/g\) and \(x=u+v\). Then
\begin{equation}\label{eq:source-coords}
 (u,v)=1,\quad ux\mid mb,\quad v\leq\delta^{-1}u,\quad
 \Om(v)\leq\Om(u),\quad u\geq L,\quad x>L,
\end{equation}
and, writing \(mb=uxd\), the \(L\)-smooth part of \(uxd\) is exactly \(m\);
in particular \(u\) is \(L\)-rough and \(x_{<L}d_{<L}=m\).
\end{lemma}

\begin{proof}
By Lemma~\ref{lem:criterion}, \(x\mid mg\), hence \(ux\mid mgu=mb\); the
remaining inequalities in \eqref{eq:source-coords} follow from \(c/b=v/u\),
the ranges of \(b,c\), and complete additivity of \(\Om\). Since \(u\mid b\)
and \(b\) is \(L\)-rough, \(u\) is \(L\)-rough. If \(u=1\) then
\(\Om(v)\leq\Om(u)=0\) forces \(v=1\) and \(c=b\), which is excluded;
therefore \(u\geq L\) and \(x=u+v>L\). Finally \(uxd=mb\) with \(b\)
\(L\)-rough and \(m\) itself \(L\)-smooth, so the \(L\)-smooth part of the
left side is \(m\).
\end{proof}

Compared with \cite[(4.3)]{DellaPietra2026} the exclusion \(x\geq L\) is
obtained here from the roughness of \(u\) rather than from a case analysis at
\(2\); this is the form that survives an arbitrary multiplier.

\begin{proposition}[source estimate]\label{prop:source-estimate}
Put \(M=A\log4\). For every fixed \(K,L\) and all sufficiently large \(N\),
\begin{equation}\label{eq:source-raw}
 \frac{T(N;m,L,A,K)}{N\rhoL}
 \ll_{A,\delta,m}
 4^K\left\{(\log L)^{-3/2}+(\log L)^{1/4-M}(\log N)^{M-3/2}\right\}.
\end{equation}
\end{proposition}

\begin{proof}
We follow \cite[Prop.~4.2]{DellaPietra2026}, indicating the four changes.

\emph{(i) The indicator.} Since \(L>m\), every prime factor of \(m\) is below
\(L\), so \(\Om_L(b,x)=\Om_L(mb,x)\). At the scale \(x\),
Lemma~\ref{lem:source-coords} and \(mb=uxd\) give
\[
 \Om_L(b,x)=\Om_L(u,x)+\Om_L(x,x)+\Om_L(d,x)
 =\Om(u)+\Om_L(x,x)+\Om_L(d,x),
\]
where the last step used \(u\leq x\) and the roughness of \(u\). Since \(x_{<L}\mid m\),
\(\Om_L(x,x)\geq\Om(x)-\Om(m)\), so regularity at \(x\) yields
\[
 \Om(u)+\Om(x)+\Om_L(d,x)\leq A\log\frac{\log x}{\log L}+K+\Om(m).
\]
Exponentiating with base \(4\) bounds the indicator by
\begin{equation}\label{eq:source-indicator}
 4^{K+\Om(m)}\left(\frac{\log x}{\log L}\right)^{M}
 4^{-\Om(u)}4^{-\Om(x)}4^{-\Om_L(d,x)} .
\end{equation}
Relative to the case \(m=2\) this costs the constant \(4^{\Om(m)-1}\).

\emph{(ii) The residual sum.} Write \(d=d_{<L}d_{\geq L}\) with
\(d_{<L}\mid m\). Since \(\Om_L\) ignores primes below \(L\),
\[
 R(D)=\sum_{\substack{d\leq D\\ d_{<L}\mid m}}4^{-\Om_L(d,x)}
 =\sum_{d_0\mid m}\ \sum_{\substack{d_1\leq D/d_0\\ \Pmin(d_1)\geq L}}
 4^{-\Om_L(d_1,x)} ,
\]
a sum of \(\tau(m)\) copies of the sum treated in
\cite[(4.6)]{DellaPietra2026}. Proposition~\ref{prop:one-variable} applies with local
factors \(1-p^{-1}\) for \(p<L\), \(1-\frac{3}{4p}+O(p^{-2})\) for
\(L\leq p\leq x\), and \(1+O(p^{-2})\) elsewhere, giving
\begin{equation}\label{eq:source-residual}
 R(D)\ll_m D\left\{(\log x)^{-3/4}(\log L)^{-1/4}+(\log(3D))^{-3/4}\right\}.
\end{equation}

\emph{(iii) The three-form sum.} Because \(\Om(v)\leq\Om(u)\),
\[
 4^{-\Om(u)}4^{-\Om(x)}\leq g_{1,L}(u)g_2(v)g_{3,L}(x),
\]
where \(g_{1,L}(n)=2^{-\Om(n)}\1_{\Pmin(n)\geq L}\),
\(g_2(n)=2^{-\Om(n)}\) and
\(g_{3,L}(n)=4^{-\Om(n)}\1_{n_{<L}\mid m}\). Setting
\(F_L(n_1,n_2,n_3)=\1_{(n_1,n_2)=1}g_{1,L}(n_1)g_2(n_2)g_{3,L}(n_3)\), the
hypothesis of Proposition~\ref{prop:dBT} holds uniformly in \(L\), and we apply it to
\(Q=(U,V,U+V)\), for which \(\rho_Q^+(p)=3p-2\) at every prime. By
\eqref{eq:axis-density} the local factor at an odd \(p\geq L\) is exactly
\[
 1+(1-p^{-1})^2\left(\frac1{2p-1}+\frac1{2p-1}+\frac1{4p-1}\right)
 =1+\frac{5}{4p}+O(p^{-2}),
\]
while at \(2<p<L\) with \(p\nmid m\) only the \(V\)-axis survives, giving
\(1+\frac1{2p}+O(p^{-2})\). At the finitely many primes dividing \(2m\) the
\(U+V\)-axis is permitted up to valuation \(v_p(m)\) and the factor is
multiplied by a bounded amount, absorbed into the constant. Mertens'
estimate and the root-sieve product \(\ll(\log X)^{-3}\) give, on
\(u\in[X,2X]\),
\begin{equation}\label{eq:source-three-form}
 \sum_{\substack{X\leq u<2X\\0<v\leq\delta^{-1}u\\(u,v)=1}}
 g_{1,L}(u)g_2(v)g_{3,L}(u+v)
 \ll_{\delta,m}X^2(\log X)^{-7/4}(\log L)^{-3/4},
\end{equation}
with every nonzero block satisfying \(X\gg_\delta L\).

\emph{(iv) The dyadic assembly.} This is identical to
\cite[Prop.~4.2]{DellaPietra2026}, with \(D=mN/(ux)\) in
\eqref{eq:source-residual}; the replacement of \(m\) for \(2\) changes \(D\)
by the bounded factor \(m/2\), hence changes \(\log(3D)\) by \(O_m(1)\).
Combining \eqref{eq:source-indicator}, \eqref{eq:source-residual} and
\eqref{eq:source-three-form} and dividing by \(N\rhoL\), one block
contributes \(\ll4^K(\log L)^{-M}(\log X)^{M-5/2}\) from the first residual
term; since \(M<3/2\) the dyadic sum over \(X\gg L\) converges to
\(\ll(\log L)^{M-3/2}\), leaving \(4^K(\log L)^{-3/2}\). The second residual
term contributes the second summand of \eqref{eq:source-raw} exactly as
there.
\end{proof}

\begin{proof}[Proof of Theorem~\ref{thm:all-k}]
For odd \(k\) the odd numbers already give \(f_k(N)\geq\lceil N/2\rceil\) by
Theorem~\ref{thm:benchmark}(1), with a far better constant than the one below;
so assume \(k\) is even. Apply the above
with \(m=k\), \(\delta=1/2\) and a fixed \(A\in(1,3/(2\log4))\). (The
argument below in fact works verbatim for every \(k\geq2\); only for even
\(k\) does it say anything.) Choose
\(K\geq K_0(A,\delta)\) from Lemma~\ref{lem:source-density}. Because
\(M=A\log4<3/2\), the first term of the right side of \eqref{eq:source-raw}
can be made at most \(\frac{1-\delta}{8}\) by fixing \(L>k\) large enough, and
with that \(L\) fixed the second term tends to \(0\), so is also at most
\(\frac{1-\delta}{8}\) for large \(N\). Hence
\(T\leq\frac{1-\delta}{4}N\rhoL\leq\frac12\abs{S}\) by
Lemma~\ref{lem:source-density}. Delete every bad \(b\) from \(S\) and call the
survivors \(B\). If distinct
\(b,d\in B\) satisfied \(b+d\mid mbd\), then relabelling so that
\(\Om(d)\leq\Om(b)\) makes \(b\) bad, a contradiction; hence \(B\) is
\(k\)-admissible. For all large \(N\),
\(\abs{B}\geq\frac12\abs{S}\geq\frac{1-\delta}{4}N\rhoL=\frac18N\rhoL\), so
\(c_k=\rhoL/8\) works.
\end{proof}

\section{The two-endpoint estimate for odd \texorpdfstring{\(k\)}{k}}
\label{sec:mixed}

From now on \(k\) is odd, and we fix the dilation
\begin{equation}\label{eq:lambda}
 \lambda=\lambda(k)=\min\{2^j:j\geq1,\ 2^j\geq k\},
\end{equation}
the least even power of two that is at least \(k\). The construction of
Theorem~\ref{thm:odd-main} is
\[
 A_N=(O_N\setminus\Gamma_N)\cup\lambda B_N\subseteq\{1,\ldots,\lambda N\},
\]
where \(O_N\) is a set of odd integers up to \(\lambda N\),
\(B_N\subseteq[\delta N,N]\) is \(k\lambda\)-admissible --- so that
\(\lambda B_N\) has no internal conflict, by
Proposition~\ref{prop:dilation} --- and \(\Gamma_N\) removes the mixed
conflicts. Two odd elements never conflict because \(k\) is odd. It therefore
remains to count the pairs \((a,b)\) with \(a\) odd, \(b\in B_N\) and
\begin{equation}\label{eq:mixed-conflict}
 a+\lambda b\mid k\lambda ab .
\end{equation}

The choice of \(\lambda\) is not cosmetic. For \(k=1\) one may take
\(\lambda=2\), as \cite{DellaPietra2026} does, but that choice is fatal as
soon as \(3\mid k\) or \(5\mid k\).

\begin{remark}[why \(\lambda=2\) fails]\label{rem:lambda-two}
Let \(n\) be odd. If \(3\mid k\) then \(n+2n=3n\) divides
\(k\cdot n\cdot 2n=2kn^2\), so \(\{n,2n\}\) is never \(k\)-admissible. With
\(\lambda=2\) each \(b\in B_N\) would then force the deletion of the odd
number \(b\) itself from \(O_N\), and that loss cancels the entire gain
\(\abs{B_N}\): the construction yields nothing. If \(5\mid k\) then
\(3n+2n=5n\) divides \(k\cdot3n\cdot2n\), so \(\{2n,3n\}\) is never
\(k\)-admissible, and the same objection applies to a positive proportion of
\(B_N\). In the coordinates of Lemma~\ref{lem:mixed-coordinates} these are
exactly the families with \(u=1\), which exist precisely when \(w+\lambda\)
divides \(k\) for some admissible \(w\). Choosing \(\lambda\geq k\) forces
\(w+\lambda>k\) and eliminates all of them; this is the sole purpose of
\eqref{eq:lambda}. These families are invisible at \(k=1\), which is why they
do not arise in \cite{DellaPietra2026}.
\end{remark}

\begin{lemma}[mixed coordinates]\label{lem:mixed-coordinates}
Let \(k\) be odd, let \(\lambda\) be as in \eqref{eq:lambda}, let \(a\) be odd
and let \(b\) be odd and \(L\)-rough, and suppose \eqref{eq:mixed-conflict}
holds. Put \(g=(a,b)\), \(w=a/g\), \(u=b/g\), \(x=w+\lambda u\),
\(\kappa=(x,k)\) and \(x'=x/\kappa\). Then \((u,w)=1\), all of \(t=g/x'\),
\(u\), \(w\), \(x\) are odd, \(x'\mid g\),
\begin{equation}\label{eq:mixed-coords}
 a=tx'w,\qquad b=tx'u,
\end{equation}
and \(\kappa\) is the \(L\)-smooth part of \(x\), so that \(\kappa\mid k\) and
\(t,u,x'\) are \(L\)-rough. Conversely every such tuple gives a conflict.

If in addition \(a\leq\lambda N\), \(b\geq\delta N\) and
\(L>\lambda(1+\delta^{-1})\), then
\begin{equation}\label{eq:u-lower}
 u\geq L,\qquad\text{and hence}\qquad x>\lambda u\geq\lambda L .
\end{equation}
\end{lemma}

\begin{proof}
Apply Lemma~\ref{lem:criterion} to the pair \((a,\lambda b)\) with multiplier
\(k\lambda\). Here \((a,\lambda b)=(a,b)=g\) because \(a\) is odd and
\(\lambda\) is a power of two; the reduced pair is \((w,\lambda u)\), coprime
because \(w\) is odd and \((w,u)=1\); and \((a+\lambda b)/g=w+\lambda u=x\),
which is odd. Lemma~\ref{lem:criterion} makes the conflict equivalent to
\(x/(x,k\lambda)\mid g\), and \(x\) is odd while \(\lambda\) is a power of
two, so \((x,k\lambda)=(x,k)=\kappa\). Hence the conflict is equivalent
to \(x'\) dividing \(g\), which gives \eqref{eq:mixed-coords} with
\(t=g/x'\). Since \(b=tx'u\) is \(L\)-rough, so are \(t\), \(x'\) and \(u\);
therefore the \(L\)-smooth part of \(x=\kappa x'\) is \(\kappa\), and
\(\kappa\mid k\) by definition. All the variables are odd because \(a\) and
\(b\) are and \(k\) is. The converse is immediate from
Lemma~\ref{lem:criterion}.

For \eqref{eq:u-lower}, suppose \(u=1\). Then \(g=b\), so
\(w=a/b\leq\lambda N/(\delta N)=\lambda\delta^{-1}\) and
\[
 x=w+\lambda\leq\lambda(1+\delta^{-1})<L .
\]
Since \(x'\mid g=b\) and \(b\) is \(L\)-rough, \(x'\) is \(L\)-rough; as
\(1\leq x'\leq x<L\), this forces \(x'=1\), that is \(x\mid k\). But
\(x=w+\lambda>\lambda\geq k\) by \eqref{eq:lambda}, and no divisor of \(k\)
exceeds \(k\). This contradiction shows \(u\neq1\). Since \(u\mid b\) and
\(b\) is \(L\)-rough, \(u\) is \(L\)-rough, so \(u\geq L\); and
\(x=w+\lambda u>\lambda u\geq\lambda L\).
\end{proof}

Fix positive \(A_b,A_o,K_b,K_o\), bases \(q_b,q_o>1\), \(\delta\in(0,1)\),
\(\lambda\) as in \eqref{eq:lambda}, and \(L>\lambda(1+\delta^{-1})\). Let
\(S\subseteq[\delta N,N]\) be the \(L\)-rough integers obeying the centered
\((A_b,K_b)\) bound, and let \(O\) be the odd \(a\leq\lambda N\) obeying the
centered \((A_o,K_o)\) bound. Put
\[
 E^{(k)}_{\rm mix}(N)=\#\{(a,b)\in O\times S:a+\lambda b\mid k\lambda ab\},
\]
\[
 \alpha=q_b^{-1},\quad\beta=q_o^{-1},\quad s=(q_bq_o)^{-1},\quad
 M=A_b\log q_b+A_o\log q_o,\quad E=M+\alpha+\beta+2s-4 .
\]

\begin{proposition}[two-endpoint estimate]\label{prop:mixed-estimate}
If \(E<-1\), then
\begin{equation}\label{eq:mixed-final}
 E^{(k)}_{\rm mix}(N)
 \ll_{k,\delta}q_b^{K_b}q_o^{K_o}\left\{N(\log L)^{-2}+o_L(N)\right\},
\end{equation}
where the normalised \(o_L\)-function may be chosen independently of
\(K_b,K_o\).
\end{proposition}

\begin{proof}
We follow \cite[Prop.~5.2]{DellaPietra2026}. Sum over the at most
\(\tau(k)\) values of \(\kappa\mid k\) separately; fix one.

\emph{Anatomy.} By Lemma~\ref{lem:mixed-coordinates}, at the scale \(x\),
\[
 \Om_L(b,x)=\Om_L(t,x)+\Om(u)+\Om(x)-\Om(\kappa),
\]
\[
 \Om_L(a,x)=\Om_L(t,x)+\Om(x)-\Om(\kappa)+\Om_L(w,x).
\]
Here we used \(u<x\), \(x'\leq x\) and the roughness of \(t,u,x'\). We also
use \(\Om_L(w,x)\geq\Om_{\geq L}(w)-1\): indeed
\(w\leq\lambda\delta^{-1}u<\lambda\delta^{-1}x\), so once
\(x\geq\lambda\delta^{-1}\) --- guaranteed by \eqref{eq:u-lower} --- the
integer \(w\) has at most one prime factor exceeding \(x\), since two would
force \(w>x^2\). (The corresponding step of \cite{DellaPietra2026} is written as an
equality; the discrepancy costs a single factor \(q_o\) and is flagged in
Section~\ref{sec:gaps}.) Exponentiating the two regularity bounds with bases
\(q_b,q_o\) majorises the joint indicator by
\begin{equation}\label{eq:mixed-indicator}
 C_k\,q_b^{K_b}q_o^{K_o}
 \left(\frac{\log x}{\log L}\right)^{M}
 \alpha^{\Om(u)}\beta^{\Om_{\geq L}(w)}s^{\Om(x)}s^{\Om_L(t,x)},
 \qquad C_k=(q_bq_o)^{\Om(k)}q_o .
\end{equation}
This is \cite[(5.4)]{DellaPietra2026} times the constant \(C_k\).

\emph{Residual sum.} With \(Y=\kappa N/(ux)\asymp_{k}N/(ux)\),
Proposition~\ref{prop:one-variable} applied to
\(h_{L,x}(t)=\1_{\Pmin(t)\geq L}s^{\Om_L(t,x)}\) gives
\begin{equation}\label{eq:mixed-residual}
 \sum_{\substack{\delta Y\leq t\leq Y\\\Pmin(t)\geq L}}s^{\Om_L(t,x)}
 \ll_{s}Y\left\{(\log x)^{s-1}(\log L)^{-s}+(\log(3Y))^{s-1}\right\},
\end{equation}
exactly as in \cite[(5.5)]{DellaPietra2026}; the bounded factor \(\kappa\)
changes \(Y\) and \(\log(3Y)\) by \(O_k(1)\).

\emph{Three-form sum.} The forms are \(Q=(U,W,\lambda U+W)\). Because
\(\lambda\) is a power of two it is invertible modulo every odd prime \(p\),
so the three root lines are distinct there and \(\rho_Q^+(p)=3p-2\), exactly
as for \(\lambda=2\); all the local computations below are therefore
independent of \(\lambda\). At \(p=2\) all three values are forced odd and the
local factor is \(1\). For \(2<p<L\) with
\(p\nmid k\) only the \(W\)-axis is permitted and \eqref{eq:axis-density}
gives the exact factor
\begin{equation}\label{eq:mixed-small-factor}
 \mathcal E_p=1+(1-p^{-1})^2\sum_{e\geq1}p^{-e}=1+\frac1p-\frac1{p^2};
\end{equation}
for the finitely many \(p\mid k\) with \(p<L\) the \(\lambda U+W\)-axis is permitted
up to valuation \(v_p(\kappa)\leq v_p(k)\), which multiplies
\eqref{eq:mixed-small-factor} by a factor between \(1\) and \(3/2\)
(certified in the supplement), contributing an \(O_k(1)\) constant. For
\(p\geq L\) the complete factor is
\begin{equation}\label{eq:mixed-large-factor}
 \mathcal E_p=1+(1-p^{-1})^2
 \left(\frac{\alpha}{p-\alpha}+\frac{\beta}{p-\beta}+\frac{s}{p-s}\right)
 =1+\frac{\alpha+\beta+s}{p}+O_{\alpha,\beta,s}(p^{-2}),
\end{equation}
with an error constant independent of \(L\). Hence, on \(u\in[X,2X]\),
\begin{equation}\label{eq:mixed-three-form}
 \sum_{\substack{X\leq u<2X\\0<w\leq\lambda\delta^{-1}u\\(u,w)=1}}
 \alpha^{\Om(u)}\1_{2\nmid w}\beta^{\Om_{\geq L}(w)}s^{\Om(\lambda u+w)}
 \1_{\Pmin(u)\geq L,\ (\lambda u+w)_{<L}\mid k}
 \ll_k X^2(\log X)^{\alpha+\beta+s-3}(\log L)^{1-\alpha-\beta-s},
\end{equation}
with \(ux\asymp_{\delta,k}X^2\). Every block with a nonzero summand has
\(X\geq L\): this is exactly \eqref{eq:u-lower}, and it is the point at which
the choice \(\lambda\geq k\) is used. Without it the blocks with \(u=1\) would
survive, and by Remark~\ref{rem:lambda-two} they carry a positive proportion
of the source whenever \(3\mid k\) or \(5\mid k\), destroying the estimate.

\emph{Assembly.} Identical to \cite[Prop.~5.2]{DellaPietra2026}. Inserting
the first term of \eqref{eq:mixed-residual} into
\eqref{eq:mixed-indicator} and using \eqref{eq:mixed-three-form}, one block
contributes \(\ll C_kq_b^{K_b}q_o^{K_o}N(\log X)^{E}(\log L)^{-E-3}\); since
\(E<-1\) the dyadic sum from \(X\geq L\) is \(\ll(\log L)^{E+1}\), leaving
exactly \(N(\log L)^{-2}\). The second residual term is \(o_L(N)\) after
normalisation by the same three-way split at \(\log X=\frac14\log N\) as
there. Summing over the \(\tau(k)\) values of \(\kappa\) multiplies the
result by \(O_k(1)\).
\end{proof}

\section{The parameter window and the proof of Theorem~\ref{thm:odd-main}}
\label{sec:assembly}

The window is that of \cite[\S6]{DellaPietra2026}; no inequality in it
mentions \(k\).

\begin{lemma}[certified window;
{\cite[\S6]{DellaPietra2026}}]\label{lem:window}
The rationals
\begin{equation}\label{eq:parameters}
\begin{aligned}
 A_b&=1.000001,&q_b&=2.48933,&A_o&=1.16312,\\
 q_o&=1.34288,&z_o&=1.34305,&c&=3.3912
\end{aligned}
\end{equation}
satisfy the five strict inequalities
\begin{equation}\label{eq:parameter-inequalities}
 A_b\log4<\tfrac32,\quad
 A_o\log z_o>z_o-1,\quad
 c\log z_o>1,\quad
 c\log q_o<1,\quad
 E<-1,
\end{equation}
with the certified one-sided margins of Table~\ref{tab:margins}.
\end{lemma}

\begin{table}[ht]
\centering
\begin{tabular}{ll}
\hline
quantity & certified one-sided margin\\
\hline
\(A_b\log4-5/2<-1\) & upper value \(<-1.1137042525\)\\
\(A_o\log z_o-(z_o-1)>0\) & \(>0.0000042730\)\\
\(c\log z_o-1>0\) & \(>0.0002111998\)\\
\(1-c\log q_o>0\) & \(>0.0002180771\)\\
\(E<-1\) & upper value \(<-1.0004077237\)\\
\hline
\end{tabular}
\caption{Strict parameter margins, recertified by rational interval
arithmetic in the supplement. Every logarithm is enclosed in a rational
interval of width below \(10^{-80}\).}
\label{tab:margins}
\end{table}

\begin{theorem}\label{thm:even-form}
Let \(k\) be odd and let \(\lambda=\lambda(k)\) be as in \eqref{eq:lambda}.
There is an absolute \(\eta_k>0\) such that for every sufficiently large
\(N\) there is a \(k\)-admissible \(A_N\subseteq\{1,\ldots,\lambda N\}\) with
\[
 \abs{A_N}\geq\frac{\lambda N}{2}+\eta_k N .
\]
\end{theorem}

\begin{proof}
Fix \eqref{eq:parameters} and \(\delta=1/2\). Choose \(z_b\in(1,2)\) with
\(A_b\log z_b>z_b-1\), and then, by Lemma~\ref{lem:source-density}, a fixed
\(K_b\) retaining a uniform positive proportion of every \(L\)-rough source
interval. This choice precedes \(L\).

For each \(N\) let \(S_N\subseteq[N/2,N]\) be the \(L\)-rough integers
satisfying the centered \((A_b,K_b)\) condition, so that
\(\abs{S_N}\geq\frac14N\rhoL\) for fixed \(L\) and all large \(N\). Delete
from \(S_N\) every \(b\) admitting \(d\leq N\), \(d\ne b\), with
\(b+d\mid k\lambda bd\) and \(\Om(d)\leq\Om(b)\), and call the survivors
\(B_N\). By the argument in the proof of Theorem~\ref{thm:all-k}, \(B_N\) is
\(k\lambda\)-admissible. Applying Proposition~\ref{prop:source-estimate} with
\(m=k\lambda\) and using \(A_b\log4<3/2\), choose
\(L>\max\{k\lambda,\ \lambda(1+\delta^{-1})\}=k\lambda\vee3\lambda\), after
\(K_b\), so large that the
first term of the right side of \eqref{eq:source-raw} is at most
\(\frac18\); the second term is then at most \(\frac18\) for large \(N\), so
the total loss is at most \(\frac14N\rhoL\leq\frac12\abs{S_N}\) and
\begin{equation}\label{eq:B-density}
 \abs{B_N}\geq\tfrac12\abs{S_N}\geq\tfrac18N\rhoL
\end{equation}
for all large \(N\).

Put \(K_o=\lceil c\log\log L\rceil\) and let \(O_N\) be the odd
\(a\leq\lambda N\) satisfying the centered \((A_o,K_o)\) condition, so that
the total number of odd \(a\leq\lambda N\) is \(\frac{\lambda N}{2}+O(1)\). By
Lemma~\ref{lem:centered-tail} and \(c\log z_o>1\),
\[
 \frac{\lambda N}{2}-\abs{O_N}\ll_\lambda Nz_o^{-K_o}
 \ll_\lambda N(\log L)^{-c\log z_o}=o_L(N\rhoL),
\]
so enlarging \(L\) gives
\begin{equation}\label{eq:odd-loss}
 \frac{\lambda N}{2}-\abs{O_N}\leq\tfrac1{32}N\rhoL .
\end{equation}
Let \(E_N=E^{(k)}_{\rm mix}(N)\) be the number of \((a,b)\in O_N\times S_N\)
with \(a+\lambda b\mid k\lambda ab\). By Proposition~\ref{prop:mixed-estimate}
and \(E<-1\),
\[
 E_N\ll_kq_b^{K_b}q_o^{K_o}\left\{N(\log L)^{-2}+o_L(N)\right\}.
\]
Since \(K_b\) is fixed and \(q_o^{K_o}\ll(\log L)^{c\log q_o}\) with
\(c\log q_o<1\), the first term is
\(\ll_kN(\log L)^{-2+c\log q_o}=o_L(N\rhoL)\), because \(\rhoL\asymp1/\log L\)
and \(-2+c\log q_o<-1\). The \(k\)-dependent constants \(C_k\) and
\(\tau(k)\), and the dilation \(\lambda=\lambda(k)\), are all fixed before
\(L\), so a further enlargement of \(L\), depending on \(k\), absorbs them. After \(L\) and hence \(K_o\) are fixed,
take \(N\) large enough for the rough source asymptotic, the second source
residual range, the odd-host bound and the normalised mixed
\(o_L\)-term. Thus
\begin{equation}\label{eq:edge-loss}
 E_N\leq\tfrac1{32}N\rhoL .
\end{equation}
The choices were made in the legal order
\[
 \text{fixed real parameters},\ \lambda(k)
 \longrightarrow K_b\longrightarrow L
 \longrightarrow K_o\longrightarrow N .
\]
Note that \(\lambda\) depends only on \(k\), so it may be fixed at the very
start; the cutoff \(L\) then depends on \(k\) through \(\lambda\), \(C_k\) and
\(\tau(k)\), and \(N\) comes last. There is no circularity.

Let \(\Gamma_N=\{a\in O_N:\exists\,b\in B_N,\ a+\lambda b\mid k\lambda ab\}\);
since \(B_N\subseteq S_N\), \(\abs{\Gamma_N}\leq E_N\). Put
\[
 A_N=(O_N\setminus\Gamma_N)\cup\{\lambda b:b\in B_N\}
 \subseteq\{1,\ldots,\lambda N\},
\]
a disjoint union because \(\lambda\) is even. Two odd members do not conflict,
since their sum is even while \(kab\) is odd. Two even members
\(\lambda b,\lambda d\) conflict exactly when \(b+d\mid k\lambda bd\), excluded
by the \(k\lambda\)-admissibility of \(B_N\) and
Proposition~\ref{prop:dilation}. Every mixed conflict was removed with
\(\Gamma_N\). Hence \(A_N\) is \(k\)-admissible, and by
\eqref{eq:B-density}, \eqref{eq:odd-loss} and \eqref{eq:edge-loss},
\[
 \abs{A_N}=\abs{O_N}-\abs{\Gamma_N}+\abs{B_N}
 \geq\frac{\lambda N}{2}-\tfrac1{32}N\rhoL-\tfrac1{32}N\rhoL
 +\tfrac18N\rhoL
 =\frac{\lambda N}{2}+\tfrac1{16}N\rhoL .
\]
The selected \(L\) is an absolute constant depending only on \(k\), so
\(\eta_k=\rhoL/16>0\).
\end{proof}

\begin{proof}[Proof of Theorem~\ref{thm:odd-main}]
Apply Theorem~\ref{thm:even-form} with \(N=\floor{X/\lambda}\). For all large
\(X\) the resulting set lies in \(\{1,\ldots,X\}\) and has at least
\[
 \frac{\lambda}{2}\floor{X/\lambda}+\eta_k\floor{X/\lambda}
 \geq\frac{X-\lambda}{2}+\frac{\eta_k}{2\lambda}X
 \geq\left(\frac12+\frac{\eta_k}{4\lambda}\right)X
\]
elements, the last step holding for \(X\) large enough that
\(\frac{\eta_k}{4\lambda}X\geq\frac\lambda2\). So
\(\varepsilon_k=\eta_k/(4\lambda)=\rhoL/(64\lambda)\) works.
\end{proof}

\section{Even multipliers: what is true and what obstructs}\label{sec:even}

For even \(k\) the only conclusion we obtain is Theorem~\ref{thm:all-k}. This
section explains why, and records the structural facts that any future
attempt will have to confront.

\subsection{No dense starting layer}

The two-layer construction of Section~\ref{sec:assembly} needs a set of density
\(1/2-o(1)\) that is \(k\)-admissible \emph{before} any deletion, because the
whole gain \(\eta_k\) is of size \(\rhoL\), far smaller than any loss
proportional to the ambient density. For odd \(k\) the odd numbers supply
that layer. By Theorem~\ref{thm:benchmark} no periodic set supplies it when \(k\) is
even, and by Proposition~\ref{prop:dilation} passing to a dilated layer \(\lambda A\)
replaces \(k\) by \(k\lambda\), which is again even. There is therefore no
scale at which a periodic layer becomes available.

\subsection{Roughness caps the source}

The output of Section~\ref{sec:source} has density \(\asymp\rhoL\), and this
is not an artefact of the write-up. What \eqref{eq:source-raw} bounds is the
bad \emph{proportion} \(T/\abs{S}\), and it bounds it by
\(\ll4^K(\log L)^{-3/2}\); any negative power of \(\log L\) would do, and the
exponent \(-3/2\) has no significance for closing the argument. What matters
is that the estimate is useful only when \(\log L\) is large, since \(4^K\)
must first be absorbed and \(K\geq K_0\) is forced by
Lemma~\ref{lem:source-density}. So \(L\) must be taken large --- and the
source \(S\) itself has density \(\frac{1-\delta}{2}\rhoL\), which tends to
\(0\) with \(L\). The method buys a small bad proportion by paying a small
absolute density, and the two cannot be decoupled: at bounded \(L\) the
estimate is vacuous. This is why the method delivers positive density and
nothing approaching \(1/2\).

\subsection{The two-adic splitting}

Corollary~\ref{cor:twoadic} does give genuine structure for even \(k\). For \(k=2\)
the conflict graph on the odd numbers is the disjoint union of its
restrictions to the classes \(1\) and \(3\) modulo \(4\). This reduces the
problem to a single class of density \(1/4\), but does not solve it: both
classes contain conflicts, the first being \((5,45)\) and \((3,15)\)
respectively, and inside a class the criterion of Lemma~\ref{lem:criterion}
becomes \(\frac{u+v}{2}\mid g\), structurally the same divisibility problem
one started with.

\subsection{A heuristic, not a theorem}

Take \(k=2\). For coprime odd \(u,v\) with \(u\equiv v\bmod4\), write
\(y=(u+v)/2\), which is odd by Corollary~\ref{cor:twoadic}. By
Lemma~\ref{lem:criterion} the pair
\[
 a=tyu,\qquad b=tyv\qquad(t\text{ odd})
\]
is a \(2\)-conflict, and \(\Om(a)\leq\Om(b)\) exactly when
\(\Om(u)\leq\Om(v)\); in that case the orientation deletes \(b\). Taking
\(u=1\) already shows that every odd multiple of \(v(v+1)/2\) with
\(v\equiv1\bmod4\) is deleted, and summing the densities
\(\asymp1/(v(u+v))\) of all these families over the pairs with
\(\Om(u)\leq\Om(v)\) gives a divergent total, since for a typical \(v\) the
inner sum \(\sum_u u^{-1}\1_{\Om(u)\leq\Om(v)}\) is unbounded. This suggests
that almost every odd number is the \(\Om\)-heavy endpoint of some
\(2\)-conflict, so that Sawin's orientation cannot retain density
\(1/2-o(1)\) from the odd numbers. We emphasise that this is a
heuristic about the \emph{method}: a divergent first moment is not a proof
that the orientation deletes almost everything, and it says nothing at all
about \(f_2\) itself, since a graph with minimum degree at least one may
still have an independent set of density \(1-o(1)\). We have not proved
either statement.

\subsection{The largest-prime-factor fibration}

The thread contains an observation of StijnC: the set of \(n\leq N\) with
\(\Pplus(n)\geq\sqrt{kN}\) has density \(\log2\approx0.693\) and, when the
largest prime factors differ, contains no conflict. The argument as posted
leaves the equal-prime case open. It can be completed, and the completed
statement explains why the observation cannot be used to beat \(1/2\).

\begin{proposition}[fibration]\label{prop:fibration}
Let \(k\geq1\), \(N\geq1\) and \(P_0>\sqrt{\max(k,2)N}\). For each prime
\(p\in(P_0,N]\) let \(A_p\subseteq\{1,\ldots,\floor{N/p}\}\) be
\(k\)-admissible. Then \(\bigcup_{P_0<p\leq N}pA_p\) is a \(k\)-admissible
subset of \(\{1,\ldots,N\}\) of size \(\sum_p\abs{A_p}\). Consequently
\[
 f_k(N)\geq\sum_{P_0<p\leq N}f_k\!\left(\floor{N/p}\right).
\]
\end{proposition}

\begin{proof}
No \(n\leq N\) is divisible by two primes exceeding \(P_0>\sqrt N\), so the
union is disjoint and each element determines its fibre.

Let \(a=pm_1\) and \(b=qm_2\) with \(p\ne q\). Then \(p\nmid b\) and
\(q\nmid a\), since otherwise \(pq\mid b\leq N<pq\). Suppose
\(kab=(a+b)h\). As \(p\mid a\) and \(p\nmid b\) we have \(p\nmid a+b\), so
\(v_p(h)=v_p(kab)\geq1\); likewise \(v_q(h)\geq1\); hence \(pq\mid h\) and
\[
 kab=(a+b)h\geq pq(a+b)>\max(k,2)N(a+b)\geq kN(a+b)>kab,
\]
using \(N\geq a\) and \(a+b>b\). This contradiction rules out conflicts
across distinct fibres.

Let \(a=pm_1\) and \(b=pm_2\) with \(m_1\ne m_2\). Put \(g'=(m_1,m_2)\) and
\(x'=(m_1+m_2)/g'\). By Lemma~\ref{lem:criterion} applied to the pair \((a,b)\),
whose gcd is \(pg'\) and whose reduced sum is \(x'\), the conflict is
equivalent to \(x'/(x',k)\mid pg'\). Now
\(x'\leq m_1+m_2\leq 2N/p<p\) because \(p>P_0>\sqrt{2N}\), so \(p\nmid x'\)
and the condition reduces to \(x'/(x',k)\mid g'\), i.e.\ to
\(m_1+m_2\mid km_1m_2\). Hence a conflict inside a fibre is exactly a
conflict in \(A_p\).
\end{proof}

\begin{remark}
Since \(\sum_{P_0<p\leq N}1/p=\log\frac{\log N}{\log P_0}+o(1)\to\log2\) for
\(P_0\asymp\sqrt N\), feeding sets of density \(\theta\) into
Proposition~\ref{prop:fibration} returns a set of density \(\theta\log2<\theta\). The
fibration transfers density downwards and can never bootstrap it; in
particular it gives no new lower bound for any \(k\).
\end{remark}

\section{Flagged items}\label{sec:gaps}

We list every point at which this manuscript relies on something it does not
prove, or at which we have altered the source presentation.

\begin{enumerate}
\item \textbf{Dependence on the companion manuscript.} Theorem~\ref{thm:odd-main}
      depends on \cite{DellaPietra2026}, which is an unrefereed proof claim.
      If its Theorem 1.1 is wrong, Theorem~\ref{thm:odd-main} is wrong for every
      odd \(k\), including \(k=1\). Theorem~\ref{thm:all-k} depends only on
      Section~\ref{sec:source}, whose method is Sawin's published one; and
      Theorem~\ref{thm:benchmark}, Lemma~\ref{lem:criterion}, Corollary~\ref{cor:twoadic},
      Proposition~\ref{prop:dilation} and Proposition~\ref{prop:fibration} are independent of both
      and are proved here in full.
\item \textbf{The two analytic inputs.} Proposition~\ref{prop:one-variable} and
      Proposition~\ref{prop:dBT} are quoted, not proved. In particular the uniformity of
      the implied constant in \eqref{eq:dBT-shape} over the multiplicative
      class, which is what allows \(F\) to depend on the cutoff \(L\), is the
      form in which \cite{DellaPietra2026} uses
      \cite[Thm.~3.1]{BretecheTenenbaum2026}; we use it identically and have
      not independently re-derived it.
\item \textbf{Ineffectivity.} Both quoted inputs carry unspecified absolute
      constants, so \(\varepsilon_k\) and \(c_k\) are not effective as
      stated. Even granting effective inputs, the cutoff \(L\) must satisfy
      \((\log L)^{1-c\log q_o}\) larger than a fixed multiple of the
      \(k\)-dependent constant \(C_k\) of \eqref{eq:mixed-indicator}, and
      \(1-c\log q_o<3\cdot10^{-4}\); the resulting \(\varepsilon_k\) is
      astronomically small. We make no claim about its size.
\item \textbf{The \(\Om_{\geq L}(w)\) majorisation.} In the proof of
      Proposition~\ref{prop:mixed-estimate} we use
      \(\Om_L(w,x)\geq\Om_{\geq L}(w)-1\) where
      \cite[Prop.~5.2]{DellaPietra2026} writes an equality. The inequality is
      correct and costs one factor \(q_o\); the equality as written is not
      justified for those \(w\) carrying a prime factor in \((x,2\delta^{-1}x]\).
      The supplement verifies exhaustively that at most one such prime factor
      can occur.
\item \textbf{Second residual ranges.} The treatment of the second residual
      term in Proposition~\ref{prop:source-estimate}(iv) and in
      Proposition~\ref{prop:mixed-estimate} is by reference to the corresponding
      passages of \cite{DellaPietra2026}, with the observation that the
      multiplier only changes the relevant \(D\) and \(Y\) by bounded
      factors. We have checked that no exponent changes, but we have not
      rewritten those two dyadic splittings in full here.
\item \textbf{The dilation.} An earlier draft of this manuscript used the
      dilation \(2\) throughout, copying \cite{DellaPietra2026}. That is
      \emph{false} for every odd \(k\) divisible by \(3\) or \(5\):
      Remark~\ref{rem:lambda-two} exhibits the offending families, which in the
      coordinates of Lemma~\ref{lem:mixed-coordinates} are exactly those with
      \(u=1\), and which carry a positive proportion of the source. The error
      was found by an adversarial audit of the draft. The fix is
      \eqref{eq:lambda}, and it is the reason
      Proposition~\ref{prop:mixed-estimate} now carries the hypothesis
      \(L>\lambda(1+\delta^{-1})\) and the conclusion \eqref{eq:u-lower}. We
      record this because the corresponding step in \cite{DellaPietra2026} ---
      the assertion that every relevant dyadic block has \(X\gg L\) --- is
      stated there without proof; it is true at \(k=1\), for the reason given
      in the proof of \eqref{eq:u-lower}, but it is not automatic.
\item \textbf{Triviality for odd \(k\).} Theorem~\ref{thm:all-k} has content
      only for even \(k\); for odd \(k\) the odd numbers already give density
      \(1/2\). We state it uniformly for convenience, not because the odd
      case is a result.
\item \textbf{Even \(k\).} We do not improve Sawin's constant at \(k=2\), we
      do not answer the thread question \(f_2(N)\geq(1/2+o(1))N\), and we
      give no explicit \(c_k\) for any \(k\geq2\).
\item \textbf{The heuristic of Section~\ref{sec:even}.} Explicitly not a theorem, as
      stated there.
\item \textbf{No formalization.} A Lean 4 development for this manuscript is
      planned and has not been written. The companion manuscript's
      formalization covers \(k=1\) and \(k=2\) only and does not certify any
      statement here.
\end{enumerate}

\section{Verification, prior work, and disclosure}\label{sec:repro}

\subsection{The exact verifier}

The supplement \path{paper/supplement/verify_certificate.py} is a standalone
standard-library Python program using exact integer and rational arithmetic;
no floating-point comparison is load-bearing. It certifies, in ten groups:
the universal periodic witness \eqref{eq:witness} and the resulting benchmark
densities of Theorem~\ref{thm:benchmark}; the conflict criterion of
Lemma~\ref{lem:criterion}; the two-adic restriction of Corollary~\ref{cor:twoadic} and the
mod-\(4\) splitting at \(k=2\); the generalised source coordinates of
Lemma~\ref{lem:source-coords}, including \(u\geq L\), \(x>L\) and the exact
\(L\)-smooth part; the mixed coordinates of Lemma~\ref{lem:mixed-coordinates} and
their converse, the exclusion \eqref{eq:u-lower} of the \(u=1\) families for
the repaired dilation together with a negative control exhibiting those
families for the dilation \(2\), and the at-most-one-large-prime fact used in
flagged item 4; the root counts \(3p-2\) for both linear-form triples and the
\(p=2\) factor; the complete small- and large-prime Euler factors for both
triples, including the modified factors at primes dividing the multiplier;
the five window inequalities of Lemma~\ref{lem:window} with rational logarithm
intervals of width below \(10^{-80}\), together with the symbolic
cancellation of the \(\log L\) exponents to \(-3/2\) in the source assembly
and \(-2\) in the mixed assembly, carried out as an identity in the formal
symbols \(M,\alpha,\beta,s\); and Proposition~\ref{prop:fibration} with its contraction
constant; and, finally, that the assembled two-layer set of
Theorem~\ref{thm:even-form} really is \(k\)-admissible in a finite box. The
recorded output is \path{paper/supplement/certificate-output.txt}.

The last group deserves a caveat. It confirms the combinatorics of the
assembly --- that the three conflict types really are all excluded --- but it
cannot establish the density gain, which is asymptotic and requires \(L\) and
\(N\) to be chosen in the legal order of Theorem~\ref{thm:even-form}. The
boxes reported there do happen to show a positive surplus over the odd
numbers, but that is a property of those particular boxes and proves nothing:
no finite computation can test the inequality of
Theorem~\ref{thm:odd-main}. What the group does test, and what it was added
to test, is that the assembled set is genuinely \(k\)-admissible --- the
property that the dilation \(2\) silently violates for \(3\mid k\) and
\(5\mid k\).

These checks certify arithmetic, local densities and the numerical window.
They are not, and are not offered as, evidence for the analytic inputs of
Section~\ref{sec:inputs}.

\subsection{Prior-work search}

A search completed on 29 July 2026 covering arXiv, the Problem 327 page and
its discussion thread, and general web search found: Sawin's paper
\cite{Sawin2026}, which treats \(k=2\) only; the upper bounds of
\cite{Leon2026} for \(k=1,2,3\); Sawin's thread comment giving the greedy
lower bound \(N/(\log N)^{0.3588\ldots+o(1)}\) for \(k=2\); and no lower
bound for \(f_k\) with \(k\geq3\) other than the odd numbers for odd \(k\),
and in particular none at all for even \(k\geq4\). We located no prior
treatment of the periodic benchmark of Theorem~\ref{thm:benchmark}. This is a dated search
report, not a categorical priority claim.

\subsection{Computational and AI-use disclosure}

AI systems were used for literature search, symbolic and numerical
exploration, adversarial auditing of the generalisation, independent
reconstruction of intermediate estimates, and drafting. The identification
of the periodic benchmark, the choice of what does and does not transfer from
\cite{DellaPietra2026}, and the decision to report Theorem~\ref{thm:odd-main} and
Theorem~\ref{thm:all-k} rather than a stronger claim are the author's. The author is
responsible for all claims and errors.

\bibliographystyle{plain}
\bibliography{references}

\end{document}
